\chapter{Discussion}

This chapter reflects on the dissertation as a whole, bringing together the research problem, design response, and evaluation findings to discuss the project’s overall contribution. It argues that consent revocation in WhatsApp-like photo sharing can be meaningfully supported through interface design, but remains partial, socially sensitive, and dependent on clear communication of its limits.


\section{Implications and Recommendations for Stakeholders}

The findings of this project carry implications beyond the prototype itself. This section outlines what the evaluation suggests for five stakeholder groups whose decisions shape how consent revocation could function in practice.

\subsection*{Users}

The evaluation showed that participants who held realistic expectations about digital persistence were more likely to trust the revocation feature than those who expected complete erasure (Theme~2). This suggests that users benefit from understanding revocable permissions as meaningful but bounded: a mechanism for withdrawing future platform-mediated access, not a guarantee that all copies have been removed. Platforms that adopt similar features should therefore invest in clear, upfront communication of limits so that users can make informed sharing decisions rather than discovering constraints after the fact.

\subsection*{Platform Designers and Developers}

Four design-level findings emerge directly from the evaluation:

\begin{enumerate}
    \item \textbf{Progressive disclosure for permissions.} The strongest usability concern (Theme~5) was that always-visible permission labels consumed excessive chat space. Designers should consider contextual menus, collapsible panels, or tap-to-reveal patterns that keep permission information accessible without dominating the conversation view.
    \item \textbf{Neutral, system-mediated revocation language.} The evaluation found that neutral wording was perceived as non-accusatory by the majority of participants (Theme~3). This validates the use of descriptive, system-like phrasing for revocation notices and should be adopted as a design guideline for any post-send permission change.
    \item \textbf{Familiar terminology for co-subject features.} The term ``co-subject'' was universally misunderstood at first encounter; four participants independently compared the concept to Instagram tagging (Theme~4). Designers implementing multi-party photo permissions should use language that maps onto existing user mental models rather than introducing unfamiliar terminology.
    \item \textbf{Permission propagation across forwarding.} The Smart Photos model's approach of linking forwarded copies to the original permission state was understood by participants and maintained sender control across redistribution. This principle should inform how real platforms handle permission inheritance for forwarded media.
\end{enumerate}

\subsection*{Platform Providers}

The evaluation demonstrates clear user demand for post-send control that goes beyond the current delete-for-everyone model. WhatsApp's recent Advanced Chat Privacy feature moves in this direction but operates at the chat level rather than the photo level \cite{whatsapp_advanced_chat_privacy}. The Smart Photos model shows what photo-level, per-item permission management could look like in practice, and the evaluation confirms that users both understood and valued this granularity. Providers considering such features should note the Theme~2 finding that some users dismissed controls they perceived as incomplete; accompanying any new permissions with honest, prominent explanations of their scope is therefore essential for maintaining user trust rather than undermining it.

\subsection*{Educators}

The evaluation revealed that participants' ability to engage with revocable permissions depended on whether they held realistic mental models of digital persistence. Several participants were surprised by the structural limits of platform-level control, such as the impossibility of preventing external-device capture. This suggests that consent education would benefit from concrete, scenario-based demonstrations of what revocation can and cannot achieve, rather than abstract policy explanations. The Smart Photos model, with its explicit permission states and visible limitation warnings, could serve as a pedagogical tool for teaching that digital consent is ongoing, partial, and context-dependent.

\subsection*{Policymakers and Regulators}

The GDPR establishes a right to withdraw consent, and the UK Online Safety Act creates obligations for platforms to reduce harm from shared content. However, the evaluation shows that technical compliance alone does not satisfy user expectations. Delete-for-everyone, for instance, operates prospectively and cannot address screenshots, saved copies, or prior viewing \cite{politou_2018}. Two policy-relevant findings emerge. First, the partial and prospective nature of revocation (Theme~2, Design Principle~2) implies that regulatory standards should set realistic expectations for what ``withdrawal of consent'' can mean in messaging contexts, rather than implying complete erasure. Second, the co-subject findings (Theme~4) suggest that multi-party consent in photo sharing is a policy area that current regulation does not adequately address. The evaluation showed that depicted individuals are widely regarded as having a legitimate stake in how shared photos are managed, yet no existing regulatory framework grants them a formal mechanism to request removal from content shared in private messaging.

\section{Limitations and Generalisability}

Several limitations should be acknowledged when interpreting the findings of this project. They fall into three categories: study design, prototype scope, and structural constraints.

(suraya bir daha bak, bir daha bak.).....
\textbf{Prototype scope.} The prototype was a front-end simulation with in-memory state, operating within a single browser window rather than across networked devices (Section~\ref{sec:tech-choices}). It did not include a backend, a database, real end-to-end encryption, or integration with WhatsApp's actual infrastructure. These constraints were deliberate and appropriate for evaluating the interaction design and permission model, but they mean that the findings do not address production-level concerns such as latency, storage, encryption key management, or cross-device synchronisation. The automated test suite used jsdom rather than a real browser engine, which further limits confidence in layout and rendering fidelity.

\textbf{Structural design constraints.} Some limitations are not specific to this prototype but inherent to any platform-level revocation mechanism. Screenshots, external-device capture, and prior viewing represent hard boundaries that no software control can address once visual content has been displayed on a screen. Five participants independently raised this concern during the evaluation. The design acknowledged this from the outset through its second design principle (expectation management over absolute control, Section~4.2.2) and through explicit pre-send and in-chat warnings, but the evaluation confirmed that this structural reality still shaped some participants' willingness to trust the feature.

Despite these limitations, the findings apply most directly to WhatsApp-like messaging contexts in which users share photos within private conversations and may later wish to revise the terms of that sharing. The underlying main principles such as dynamic consent, multi-party awareness, expectation management, and neutral communication of access changes, are not specific to WhatsApp and could inform the design of revocable sharing features on other platforms where users exchange personal media.

\section{Future Directions}

The limitations and evaluation findings identified above point to several directions for future research, each grounded in a specific gap that the current project could not address.

\textbf{Live-prototype evaluation.} The present study used screenshot-based interviews, which captured participants' interpretive reasoning but not their real-time interaction behaviour. A follow-up study in which participants use the running prototype to send, revoke, and forward photos across simulated conversations would strengthen ecological validity and reveal usability issues tied to timing, responsiveness, and sequential decision-making that static materials cannot capture.

\textbf{Broader and more diverse samples.} The current sample was not gender-balanced. Evaluating the Smart Photos model with gender-balanced recruitment, and non-student populations would test whether the findings specifically around the social meaning of revocation (Theme~3) and the acceptability of co-subject governance models (Theme~4), generalise beyond the present demographic.

\textbf{Co-subject identification and governance.} The evaluation surfaced an unresolved tension between manual tagging, which depends on sender willingness, and automated detection, which introduces accuracy and privacy risks. Future work could explore hybrid approaches in which lightweight detection prompts the sender to confirm co-subject presence, reducing friction while preserving sender agency. The governance question about whether the co-subjects should hold equal, superior, or veto-level control also warrants dedicated experimental comparison.

\textbf{Integration with real platform infrastructure.} Testing the permission model within an actual messaging platform's architecture would reveal implementation constraints that a front-end simulation cannot surface, including latency, storage overhead, encryption compatibility, and the behaviour of permission propagation under unreliable network conditions.

\textbf{Longitudinal adoption.} The evaluation captured first-exposure reactions. Diary studies or longitudinal field trials tracking how users engage with revocable permissions over weeks or months would address whether the features retain their perceived value as routine settles, or whether users simplify their behaviour by defaulting to a small subset of permission configurations.



